\begin{question}{Representation Surgery}

\begin{subquestion}{\textbf{Linear vs. Non-Linear Classifier}}

No. While $X$ linearly guarding $Z$ implies that no affine (linear) predictor can outperform the trivially attainable loss $L_{\tau}$, it does not restrict non-linear classifiers.
\end{subquestion}

\begin{subquestion}{\textbf{Proof of Linear Guardedness}}

To prove that $X$ linearly guards $Z$, we must show that for any affine predictor $f(x) = Ax + b$, the expected loss satisfies $\mathbb{E}_{X,Z}[\mathcal{L}(f(X), Z)] \ge L_{\tau}$.

By the law of total expectation, we decompose the global loss into class-conditional expectations:
$$\mathbb{E}_{X,Z}[\mathcal{L}(AX + b, Z)] = \sum_{i=1}^{k} \mathbb{P}(Z=i) \mathbb{E}_{X|Z=i} [\mathcal{L}(AX + b, Z=i)]$$

The cross-entropy loss function $\mathcal{L}(y, Z)$ is convex with respect to its first argument. Applying Jensen's Inequality to each class-conditional expectation yields:
$$\mathbb{E}_{X|Z=i} [\mathcal{L}(AX + b, Z=i)] \ge \mathcal{L}(\mathbb{E}_{X|Z=i}[AX + b], Z=i)$$

Using the linearity of expectation and substituting the given condition $\mathbb{E}[X|Z=i] = \mathbb{E}[X]$:
$$\mathcal{L}(\mathbb{E}_{X|Z=i}[AX + b], Z=i) = \mathcal{L}(A\mathbb{E}[X|Z=i] + b, Z=i) = \mathcal{L}(A\mathbb{E}[X] + b, Z=i)$$

Let $b' = A\mathbb{E}[X] + b \in \mathbb{R}^k$. Note that $b'$ is a constant vector that does not depend on the specific class index $i$. Recombining the terms into the global expectation:
$$\sum_{i=1}^{k} \mathbb{P}(Z=i) \mathcal{L}(b', Z=i) = \mathbb{E}_Z[\mathcal{L}(b', Z)]$$

By Definition 1, $L_{\tau}$ is the infimum over all constant predictors $b \in \mathbb{R}^k$. Therefore:
$$\mathbb{E}_Z[\mathcal{L}(b', Z)] \ge \inf_{b \in \mathbb{R}^k} \mathbb{E}_Z [\mathcal{L}(b, Z)] = L_{\tau}$$

This confirms that no affine classifier can achieve a loss lower than the trivially attainable loss, thus $X$ linearly guards $Z$.
\end{subquestion}

\begin{subquestion}{\textbf{Gradient of Cross-Entropy Loss}}

\subsubsection*{Boundedness of Partial Derivatives}

To prove that the partial derivatives are bounded, we first calculate the gradient of the loss $\mathcal{L}(y, Z=i)$ with respect to each component $y_j$ of the logit vector $y$.

Using the form $\mathcal{L}(y, Z=i) = -y_i + \ln\left(\sum_{m=1}^{k} \exp(y_m)\right)$, we compute the partial derivative $\frac{\partial \mathcal{L}}{\partial y_j}$:
$$\frac{\partial \mathcal{L}}{\partial y_j} = \frac{\partial}{\partial y_j} \left[ -y_i + \ln\left(\sum_{m=1}^{k} \exp(y_m)\right) \right]$$

Applying the chain rule to the logarithm term:
$$\frac{\partial \mathcal{L}}{\partial y_j} = -\mathbb{1}\{j=i\} + \frac{1}{\sum_{m=1}^{k} \exp(y_m)} \cdot \exp(y_j)$$
$$\frac{\partial \mathcal{L}}{\partial y_j} = \text{softmax}(y)_j - \mathbb{1}\{j=i\}$$

Where $\mathbb{1}\{j=i\}$ is the indicator function (1 if $j=i$, 0 otherwise). 

By the definition of the softmax function, $0 < \text{softmax}(y)_j < 1$ for all $j \in \{1, \dots, k\}$. Since the indicator function is either $0$ or $1$, the partial derivative is bounded:
$$-1 < \frac{\partial \mathcal{L}}{\partial y_j} < 1$$
Thus, all partial derivatives of the cross-entropy loss with respect to $y$ are bounded.

\subsubsection*{Symmetric Partial Derivatives for Incorrect Classes}

We want to show that for any $j_1, j_2 \neq i$, the partial derivative of the loss with respect to $y_i$ is identical. 

Using the gradient derived in the previous sub-question:
$$\frac{\partial \mathcal{L}(y, Z=k)}{\partial y_i} = \text{softmax}(y)_i - \mathbb{1}\{i=k\}$$

Let us evaluate this for the two cases where the ground truth concept $Z$ is $j_1$ and $j_2$:

1. For $Z = j_1$: 
Since $j_1 \neq i$ by assumption, the indicator function $\mathbb{1}\{i=j_1\} = 0$. Thus:
$$\frac{\partial}{\partial y_i} \mathcal{L}(y, Z=j_1) = \text{softmax}(y)_i - 0 = \text{softmax}(y)_i$$

2. For $Z = j_2$: 
Similarly, since $j_2 \neq i$ by assumption, the indicator function $\mathbb{1}\{i=j_2\} = 0$. Thus:
$$\frac{\partial}{\partial y_i} \mathcal{L}(y, Z=j_2) = \text{softmax}(y)_i - 0 = \text{softmax}(y)_i$$

Since both expressions evaluate to $\text{softmax}(y)_i$, we have:
$$\frac{\partial}{\partial y_i} \mathcal{L}(y, Z=j_1) = \frac{\partial}{\partial y_i} \mathcal{L}(y, Z=j_2)$$

Furthermore, because the softmax output $\text{softmax}(y)_i = \frac{\exp(y_i)}{\sum \exp(y_m)}$ is always strictly positive, this common derivative is non-zero ($\ne 0$), as required.

\end{subquestion}

\begin{subquestion}
Skipped.
\end{subquestion}

\begin{subquestion}{\textbf{Equivalence of Mean Equality and Zero Cross-Covariance}}

We aim to prove that $\mathbb{E}[X|Z=i] = \mathbb{E}[X]$ for all $i \in \{1, \dots, k\}$ if and only if $\Sigma_{X,\tilde{Z}} = 0$, given $\mathbb{P}(Z=i) \neq 0$.

\subsubsection*{Forward Direction: $\mathbb{E}[X|Z=i] = \mathbb{E}[X] \implies \Sigma_{X,\tilde{Z}} = 0$}
Recall the computational formula for cross-covariance:
$$\Sigma_{X,\tilde{Z}} = \mathbb{E}[X\tilde{Z}^\top] - \mathbb{E}[X]\mathbb{E}[\tilde{Z}]^\top$$
Using the law of total expectation on the first term:
$$\mathbb{E}[X\tilde{Z}^\top] = \sum_{i=1}^k \mathbb{P}(Z=i) \mathbb{E}[X\tilde{Z}^\top | Z=i]$$
Given that $\tilde{Z}$ is a one-hot encoding, $\tilde{Z} | Z=i$ is the constant vector $e_i$. Thus:
$$\mathbb{E}[X\tilde{Z}^\top] = \sum_{i=1}^k \mathbb{P}(Z=i) \mathbb{E}[X | Z=i] e_i^\top$$
Substituting the assumption $\mathbb{E}[X|Z=i] = \mathbb{E}[X]$:
$$\mathbb{E}[X\tilde{Z}^\top] = \mathbb{E}[X] \sum_{i=1}^k \mathbb{P}(Z=i) e_i^\top = \mathbb{E}[X] \mathbb{E}[\tilde{Z}]^\top$$
Therefore, $\Sigma_{X,\tilde{Z}} = \mathbb{E}[X]\mathbb{E}[\tilde{Z}]^\top - \mathbb{E}[X]\mathbb{E}[\tilde{Z}]^\top = 0$.

\subsubsection*{Backward Direction: $\Sigma_{X,\tilde{Z}} = 0 \implies \mathbb{E}[X|Z=i] = \mathbb{E}[X]$}
If $\Sigma_{X,\tilde{Z}} = 0$, then for each class $i$ (corresponding to the $i^{th}$ column of the matrix):
$$\mathbb{E}[X\tilde{Z}_i] - \mathbb{E}[X]\mathbb{E}[\tilde{Z}_i] = 0$$
Since $\tilde{Z}_i$ is an indicator variable $\mathbb{1}\{Z=i\}$, we have $\mathbb{E}[\tilde{Z}_i] = \mathbb{P}(Z=i)$ and $\mathbb{E}[X\tilde{Z}_i] = \mathbb{P}(Z=i)\mathbb{E}[X|Z=i]$. Substituting these:
$$\mathbb{P}(Z=i)\mathbb{E}[X|Z=i] - \mathbb{E}[X]\mathbb{P}(Z=i) = 0$$
Factoring out $\mathbb{P}(Z=i)$:
$$\mathbb{P}(Z=i) \left( \mathbb{E}[X|Z=i] - \mathbb{E}[X] \right) = 0$$
Since $\mathbb{P}(Z=i) \neq 0$, it must hold that $\mathbb{E}[X|Z=i] = \mathbb{E}[X]$ for all $i$.
\end{subquestion}

\begin{subquestion}{\textbf{Null Space Condition for Linear Guarding}}

A representation $r(X)$ linearly guards $Z$ if and only if the cross-covariance $\text{Cov}(r(X), \tilde{Z}) = 0$. Using the affine map $r(X) = PX + b$, we apply the linearity of covariance:
$$\text{Cov}(PX + b, \tilde{Z}) = \text{Cov}(PX, \tilde{Z}) + \text{Cov}(b, \tilde{Z})$$
Since $b$ is a constant, $\text{Cov}(b, \tilde{Z}) = 0$. By the properties of the covariance operator for matrices:
$$\text{Cov}(PX, \tilde{Z}) = P \text{Cov}(X, \tilde{Z}) = P \Sigma_{X,\tilde{Z}}$$
Therefore, $r(X)$ linearly guards $Z$ if and only if $P \Sigma_{X,\tilde{Z}} = 0$.
\end{subquestion}

\begin{subquestion}

\subsubsection*{Direction $(\implies):$ $r(\mathbf{X})$ linearly guards $Z \implies \mathbf{P} \boldsymbol{\Sigma}_{\mathbf{X}, \tilde{Z}} = \mathbf{0}$}

Assume that the representation $r(\mathbf{X}) = \mathbf{P}\mathbf{X} + \mathbf{b}$ linearly guards the concept $Z$. By definition of linear guarding, the cross-covariance between the transformed representation and the target concept must vanish:
$$\boldsymbol{\Sigma}_{r(\mathbf{X}), \tilde{Z}} = \mathbf{0}$$

Substituting the expression for $r(\mathbf{X})$, we have:
$$Cov(\mathbf{P}\mathbf{X} + \mathbf{b}, \tilde{Z}) = \mathbf{0}$$

We now apply the \textbf{linearity of covariance}. Since covariance is invariant to translations by a constant vector $\mathbf{b}$, and linear with respect to the matrix $\mathbf{P}$, we can simplify the expression as follows:
$$Cov(\mathbf{P}\mathbf{X} + \mathbf{b}, \tilde{Z}) = Cov(\mathbf{P}\mathbf{X}, \tilde{Z}) = \mathbf{P} Cov(\mathbf{X}, \tilde{Z})$$

Substituting the definition of the original cross-covariance $\boldsymbol{\Sigma}_{\mathbf{X}, \tilde{Z}} = Cov(\mathbf{X}, \tilde{Z})$, we arrive at:
$$\mathbf{P} \boldsymbol{\Sigma}_{\mathbf{X}, \tilde{Z}} = \mathbf{0}$$

This result confirms that if $r(\mathbf{X})$ is guarded, the matrix $\mathbf{P}$ must map the columns of the covariance matrix to the zero vector.

\subsubsection*{Direction $(\impliedby):$ $\mathbf{P} \boldsymbol{\Sigma}_{\mathbf{X}, \tilde{Z}} = \mathbf{0} \implies r(\mathbf{X})$ linearly guards $Z$}

Now, we assume that the algebraic condition $\mathbf{P} \boldsymbol{\Sigma}_{\mathbf{X}, \tilde{Z}} = \mathbf{0}$ holds. We must show that this condition is sufficient to guarantee that the representation $r(\mathbf{X}) = \mathbf{P}\mathbf{X} + \mathbf{b}$ linearly guards $Z$.

We examine the cross-covariance of the transformed representation:
$$\boldsymbol{\Sigma}_{r(\mathbf{X}), \tilde{Z}} = Cov(\mathbf{P}\mathbf{X} + \mathbf{b}, \tilde{Z})$$

Using the \textbf{linearity of covariance}, we apply translation invariance to the constant vector $\mathbf{b}$ and homogeneity to the matrix $\mathbf{P}$:
$$\boldsymbol{\Sigma}_{r(\mathbf{X}), \tilde{Z}} = \mathbf{P} Cov(\mathbf{X}, \tilde{Z})$$

Substituting the definition of the original cross-covariance matrix $\boldsymbol{\Sigma}_{\mathbf{X}, \tilde{Z}}$:
$$\boldsymbol{\Sigma}_{r(\mathbf{X}), \tilde{Z}} = \mathbf{P} \boldsymbol{\Sigma}_{\mathbf{X}, \tilde{Z}}$$

By our initial assumption, $\mathbf{P} \boldsymbol{\Sigma}_{\mathbf{X}, \tilde{Z}} = \mathbf{0}$. Therefore:
$$\boldsymbol{\Sigma}_{r(\mathbf{X}), \tilde{Z}} = \mathbf{0}$$

Since the cross-covariance between $r(\mathbf{X})$ and $\tilde{Z}$ is the zero matrix, it follows from the established equivalence between zero covariance and linear guarding that $r(\mathbf{X})$ linearly guards $Z$.
\end{subquestion}

\begin{subquestion}

\subsubsection*{(i)}

\begin{align*}
D(\mathbf{P}, \mathbf{b}) &= \mathbb{E}[|\mathbf{P}\mathbf{X} - \mathbf{X} + \mathbf{b}|^2] \\
&= \mathbb{E}[|\mathbf{P}\mathbf{X} - \mathbf{X}|^2 + |\mathbf{b}|^2 + 2(\mathbf{P}\mathbf{X} - \mathbf{X} | \mathbf{b})] \\
&= \mathbb{E}[|\mathbf{P}\mathbf{X} - \mathbf{X}|^2] + |\mathbf{b}|^2 + 2\mathbb{E}[(\mathbf{P}\mathbf{X} - \mathbf{X} | \mathbf{b})] \\
&= \mathbb{E}[|\mathbf{P}\mathbf{X} - \mathbf{X}|^2] + |\mathbf{b}|^2 + 2(\mathbb{E}[(\mathbf{P} - \mathbf{I})\mathbf{X}] | \mathbf{b}) \\
&= \mathbb{E}[|\mathbf{P}\mathbf{X} - \mathbf{X}|^2] + |\mathbf{b}|^2 + 2((\mathbf{P} - \mathbf{I})\mathbb{E}[\mathbf{X}] | \mathbf{b}) \\
\text{Since } \mathbb{E}[\mathbf{X}] &= \mathbf{0}: \\
&= \mathbb{E}[|\mathbf{P}\mathbf{X} - \mathbf{X}|^2] + |\mathbf{b}|^2 + 2(\mathbf{0} | \mathbf{b}) \\
&= \mathbb{E}[|\mathbf{P}\mathbf{X} - \mathbf{X}|^2] + |\mathbf{b}|^2
\end{align*}

\subsubsection*{(ii)}

Since $\mathbf{b}$ only adds to the distance the optional value is $\mathbf{0}$.

\end{subquestion}

\begin{subquestion}
\subsubsection*{Proof that $\Pstar$ satisfies Eq. (7)}
\begin{align*}
           \Pstar &= \I - \SigXZ \SigXZ^{+} \\
    \Pstar \SigXZ &= (\I - \SigXZ \SigXZ^{+}) \SigXZ \\
                  &= \I \SigXZ - (\SigXZ \SigXZ^{+} \SigXZ) \\
                  &= \SigXZ - \SigXZ \justification{$A A^+ A = A$} \\
    \Pstar \SigXZ &= \mathbf{0}
\end{align*}
\subsubsection*{Proof that $\Pstar$ minimizes Eq. (8)}
\end{subquestion}

\begin{subquestion}

\end{subquestion}

\end{question}